% Paste-ready Results section for the IEEE draft.
%
% Add this package in the preamble if it is not already present:
% \usepackage{booktabs}
%
% Add this line after \begin{document}, before this section:
% % AUTO-GENERATED -- DO NOT EDIT BY HAND
% Generated from paper/numbers.json
\providecommand{\xspace}{}

% numbers.json:ablation_gate_output_cosine
\newcommand{\AblationGateOutputCosine}{0.996\xspace}
% numbers.json:ablation_k_output_cosine
\newcommand{\AblationKOutputCosine}{1\xspace}
% numbers.json:ablation_q_output_cosine
\newcommand{\AblationQOutputCosine}{0.993\xspace}
% numbers.json:ablation_state_output_cosine
\newcommand{\AblationStateOutputCosine}{0.994\xspace}
% numbers.json:ablation_v_output_cosine
\newcommand{\AblationVOutputCosine}{1\xspace}
% numbers.json:h100_baseline_energy_mj
\newcommand{\HOneHundredBaselineEnergyMj}{99.8\,mJ\xspace}
% numbers.json:h100_baseline_latency_us
\newcommand{\HOneHundredBaselineLatencyUs}{285\,$\mu$s\xspace}
% numbers.json:h100_baseline_power_w
\newcommand{\HOneHundredBaselinePowerW}{350\,W\xspace}
% numbers.json:headline_energy_efficiency_vs_h100
\newcommand{\HeadlineEnergyEfficiencyVsHOneHundred}{216.28$\times$\xspace}
% numbers.json:headline_energy_efficiency_vs_usc
\newcommand{\HeadlineEnergyEfficiencyVsUsc}{20.588$\times$\xspace}
% numbers.json:headline_speedup_vs_h100
\newcommand{\HeadlineSpeedupVsHOneHundred}{5.544$\times$\xspace}
% numbers.json:headline_speedup_vs_usc
\newcommand{\HeadlineSpeedupVsUsc}{1.229$\times$\xspace}
% numbers.json:hls_csim_freshness_status
\newcommand{\HlsCsimFreshnessStatus}{current\xspace}
% numbers.json:hls_csim_vectors_passed
\newcommand{\HlsCsimVectorsPassed}{64\xspace}
% numbers.json:hls_csim_vectors_total
\newcommand{\HlsCsimVectorsTotal}{64\xspace}
% numbers.json:offchip_mxfp4_b32_naive_three_pass_bytes_per_token
\newcommand{\OffchipMxfpFourBThirtyTwoNaiveThreePassBytesPerToken}{835584\xspace}
% numbers.json:offchip_mxfp4_b32_persistent_bytes_per_token
\newcommand{\OffchipMxfpFourBThirtyTwoPersistentBytesPerToken}{0\xspace}
% numbers.json:offchip_mxfp4_b32_state_bytes
\newcommand{\OffchipMxfpFourBThirtyTwoStateBytes}{278528\xspace}
% numbers.json:ours_mxfp4_b16_pk16_pv8_bram_pct
\newcommand{\OursMxfpFourBSixteenPkSixteenPvEightBramPct}{0.37\%\xspace}
% numbers.json:ours_mxfp4_b16_pk16_pv8_dsp_pct
\newcommand{\OursMxfpFourBSixteenPkSixteenPvEightDspPct}{7.27\%\xspace}
% numbers.json:ours_mxfp4_b16_pk16_pv8_latency_us
\newcommand{\OursMxfpFourBSixteenPkSixteenPvEightLatencyUs}{9868.144\,$\mu$s\xspace}
% numbers.json:ours_mxfp4_b16_pk16_pv8_lut_pct
\newcommand{\OursMxfpFourBSixteenPkSixteenPvEightLutPct}{6.56\%\xspace}
% numbers.json:ours_mxfp4_b16_pk16_pv8_power_w
\newcommand{\OursMxfpFourBSixteenPkSixteenPvEightPowerW}{4.947\,W\xspace}
% numbers.json:ours_mxfp4_b16_pk32_pv16_bram_pct
\newcommand{\OursMxfpFourBSixteenPkThirtyTwoPvSixteenBramPct}{0.57\%\xspace}
% numbers.json:ours_mxfp4_b16_pk32_pv16_dsp_pct
\newcommand{\OursMxfpFourBSixteenPkThirtyTwoPvSixteenDspPct}{28.72\%\xspace}
% numbers.json:ours_mxfp4_b16_pk32_pv16_latency_us
\newcommand{\OursMxfpFourBSixteenPkThirtyTwoPvSixteenLatencyUs}{8240.176\,$\mu$s\xspace}
% numbers.json:ours_mxfp4_b16_pk32_pv16_lut_pct
\newcommand{\OursMxfpFourBSixteenPkThirtyTwoPvSixteenLutPct}{24.96\%\xspace}
% numbers.json:ours_mxfp4_b16_pk32_pv16_power_w
\newcommand{\OursMxfpFourBSixteenPkThirtyTwoPvSixteenPowerW}{11.004\,W\xspace}
% numbers.json:ours_mxfp4_b32_bram_pct
\newcommand{\OursMxfpFourBThirtyTwoBramPct}{3.67\%\xspace}
% numbers.json:ours_mxfp4_b32_csynth_fmax_mhz
\newcommand{\OursMxfpFourBThirtyTwoCsynthFmaxMhz}{342.47\,MHz\xspace}
% numbers.json:ours_mxfp4_b32_csynth_latency_cycles
\newcommand{\OursMxfpFourBThirtyTwoCsynthLatencyCycles}{12852\,cycles\xspace}
% numbers.json:ours_mxfp4_b32_dsp_pct
\newcommand{\OursMxfpFourBThirtyTwoDspPct}{7.27\%\xspace}
% numbers.json:ours_mxfp4_b32_energy_mj
\newcommand{\OursMxfpFourBThirtyTwoEnergyMj}{0.461\,mJ\xspace}
% numbers.json:ours_mxfp4_b32_impl_wns_ns
\newcommand{\OursMxfpFourBThirtyTwoImplWnsNs}{0.065\,ns\xspace}
% numbers.json:ours_mxfp4_b32_latency_cycles
\newcommand{\OursMxfpFourBThirtyTwoLatencyCycles}{12852\,cycles\xspace}
% numbers.json:ours_mxfp4_b32_latency_source_kind
\newcommand{\OursMxfpFourBThirtyTwoLatencySourceKind}{csynth\_estimate\xspace}
% numbers.json:ours_mxfp4_b32_latency_us
\newcommand{\OursMxfpFourBThirtyTwoLatencyUs}{51.408\,$\mu$s\xspace}
% numbers.json:ours_mxfp4_b32_lut_pct
\newcommand{\OursMxfpFourBThirtyTwoLutPct}{6.64\%\xspace}
% numbers.json:ours_mxfp4_b32_power_w
\newcommand{\OursMxfpFourBThirtyTwoPowerW}{8.976\,W\xspace}
% numbers.json:ours_mxfp4_b32_synth_wns_ns
\newcommand{\OursMxfpFourBThirtyTwoSynthWnsNs}{1.118\,ns\xspace}
% numbers.json:ours_mxfp4_b32_tokens_cosim
\newcommand{\OursMxfpFourBThirtyTwoTokensCosim}{128\,tokens\xspace}
% numbers.json:ours_mxfp4_pk32_pv16_b32_bram_pct
\newcommand{\OursMxfpFourPkThirtyTwoPvSixteenBThirtyTwoBramPct}{0.57\%\xspace}
% numbers.json:ours_mxfp4_pk32_pv16_b32_dsp_pct
\newcommand{\OursMxfpFourPkThirtyTwoPvSixteenBThirtyTwoDspPct}{28.72\%\xspace}
% numbers.json:ours_mxfp4_pk32_pv16_b32_latency_us
\newcommand{\OursMxfpFourPkThirtyTwoPvSixteenBThirtyTwoLatencyUs}{8240.176\,$\mu$s\xspace}
% numbers.json:ours_mxfp4_pk32_pv16_b32_lut_pct
\newcommand{\OursMxfpFourPkThirtyTwoPvSixteenBThirtyTwoLutPct}{24.96\%\xspace}
% numbers.json:ours_mxfp4_pk32_pv16_b32_power_w
\newcommand{\OursMxfpFourPkThirtyTwoPvSixteenBThirtyTwoPowerW}{11.004\,W\xspace}
% numbers.json:ours_mxfp4_pk8_pv4_b32_bram_pct
\newcommand{\OursMxfpFourPkEightPvFourBThirtyTwoBramPct}{0.47\%\xspace}
% numbers.json:ours_mxfp4_pk8_pv4_b32_dsp_pct
\newcommand{\OursMxfpFourPkEightPvFourBThirtyTwoDspPct}{1.86\%\xspace}
% numbers.json:ours_mxfp4_pk8_pv4_b32_latency_us
\newcommand{\OursMxfpFourPkEightPvFourBThirtyTwoLatencyUs}{13034.064\,$\mu$s\xspace}
% numbers.json:ours_mxfp4_pk8_pv4_b32_lut_pct
\newcommand{\OursMxfpFourPkEightPvFourBThirtyTwoLutPct}{2.04\%\xspace}
% numbers.json:ours_mxfp4_pk8_pv4_b32_power_w
\newcommand{\OursMxfpFourPkEightPvFourBThirtyTwoPowerW}{3.86\,W\xspace}
% numbers.json:qwen_capture_status
\newcommand{\QwenCaptureStatus}{available\xspace}
% numbers.json:state_storage_mxfp4_b32_bytes
\newcommand{\StateStorageMxfpFourBThirtyTwoBytes}{278528\xspace}
% numbers.json:state_storage_mxfp4_b32_reduction_vs_bf16_pct
\newcommand{\StateStorageMxfpFourBThirtyTwoReductionVsBfSixteenPct}{73.438\%\xspace}
% numbers.json:state_storage_mxfp4_b32_reduction_vs_fp32_pct
\newcommand{\StateStorageMxfpFourBThirtyTwoReductionVsFpThirtyTwoPct}{86.719\%\xspace}
% numbers.json:synthetic_boundary_state_mxfp4_b16_checkpoint_token
\newcommand{\SyntheticBoundaryStateMxfpFourBSixteenCheckpointToken}{256\,tokens\xspace}
% numbers.json:synthetic_boundary_state_mxfp4_b16_output_cosine
\newcommand{\SyntheticBoundaryStateMxfpFourBSixteenOutputCosine}{0.747\xspace}
% numbers.json:synthetic_boundary_state_mxfp4_b16_state_rel_l2
\newcommand{\SyntheticBoundaryStateMxfpFourBSixteenStateRelLTwo}{0.647\xspace}
% numbers.json:synthetic_boundary_state_mxfp4_b32_checkpoint_token
\newcommand{\SyntheticBoundaryStateMxfpFourBThirtyTwoCheckpointToken}{256\,tokens\xspace}
% numbers.json:synthetic_boundary_state_mxfp4_b32_output_cosine
\newcommand{\SyntheticBoundaryStateMxfpFourBThirtyTwoOutputCosine}{0.735\xspace}
% numbers.json:synthetic_boundary_state_mxfp4_b32_state_rel_l2
\newcommand{\SyntheticBoundaryStateMxfpFourBThirtyTwoStateRelLTwo}{0.671\xspace}
% numbers.json:synthetic_boundary_state_mxfp8_b16_checkpoint_token
\newcommand{\SyntheticBoundaryStateMxfpEightBSixteenCheckpointToken}{256\,tokens\xspace}
% numbers.json:synthetic_boundary_state_mxfp8_b16_output_cosine
\newcommand{\SyntheticBoundaryStateMxfpEightBSixteenOutputCosine}{0.945\xspace}
% numbers.json:synthetic_boundary_state_mxfp8_b16_state_rel_l2
\newcommand{\SyntheticBoundaryStateMxfpEightBSixteenStateRelLTwo}{0.282\xspace}
% numbers.json:synthetic_boundary_state_mxfp8_b32_checkpoint_token
\newcommand{\SyntheticBoundaryStateMxfpEightBThirtyTwoCheckpointToken}{256\,tokens\xspace}
% numbers.json:synthetic_boundary_state_mxfp8_b32_output_cosine
\newcommand{\SyntheticBoundaryStateMxfpEightBThirtyTwoOutputCosine}{0.945\xspace}
% numbers.json:synthetic_boundary_state_mxfp8_b32_state_rel_l2
\newcommand{\SyntheticBoundaryStateMxfpEightBThirtyTwoStateRelLTwo}{0.282\xspace}
% numbers.json:synthetic_drift_mxfp4_final_output_cosine
\newcommand{\SyntheticDriftMxfpFourFinalOutputCosine}{0.643\xspace}
% numbers.json:synthetic_drift_mxfp4_final_state_rel_l2
\newcommand{\SyntheticDriftMxfpFourFinalStateRelLTwo}{0.865\xspace}
% numbers.json:synthetic_drift_mxfp8_final_output_cosine
\newcommand{\SyntheticDriftMxfpEightFinalOutputCosine}{0.933\xspace}
% numbers.json:synthetic_drift_mxfp8_final_state_rel_l2
\newcommand{\SyntheticDriftMxfpEightFinalStateRelLTwo}{0.344\xspace}
% numbers.json:synthetic_drift_tokens
\newcommand{\SyntheticDriftTokens}{1024\,tokens\xspace}
% numbers.json:synthetic_int4_output_cosine
\newcommand{\SyntheticIntFourOutputCosine}{0.967\xspace}
% numbers.json:synthetic_mxfp4_b16_output_cosine
\newcommand{\SyntheticMxfpFourBSixteenOutputCosine}{0.982\xspace}
% numbers.json:synthetic_mxfp4_b16_output_rel_l2
\newcommand{\SyntheticMxfpFourBSixteenOutputRelLTwo}{0.189\xspace}
% numbers.json:synthetic_mxfp4_b16_state_mxfp8_output_cosine
\newcommand{\SyntheticMxfpFourBSixteenStateMxfpEightOutputCosine}{0.988\xspace}
% numbers.json:synthetic_mxfp4_b16_state_rel_l2
\newcommand{\SyntheticMxfpFourBSixteenStateRelLTwo}{0.117\xspace}
% numbers.json:synthetic_mxfp4_b32_output_cosine
\newcommand{\SyntheticMxfpFourBThirtyTwoOutputCosine}{0.982\xspace}
% numbers.json:synthetic_mxfp4_b32_output_rel_l2
\newcommand{\SyntheticMxfpFourBThirtyTwoOutputRelLTwo}{0.191\xspace}
% numbers.json:synthetic_mxfp4_b32_state_mxfp8_output_cosine
\newcommand{\SyntheticMxfpFourBThirtyTwoStateMxfpEightOutputCosine}{0.988\xspace}
% numbers.json:synthetic_mxfp4_b32_state_rel_l2
\newcommand{\SyntheticMxfpFourBThirtyTwoStateRelLTwo}{0.118\xspace}
% numbers.json:synthetic_stress_combined_stress_state_mxfp4_b16_mean_output_cosine
\newcommand{\SyntheticStressCombinedStressStateMxfpFourBSixteenMeanOutputCosine}{0.974\xspace}
% numbers.json:synthetic_stress_combined_stress_state_mxfp4_b16_worst_output_cosine
\newcommand{\SyntheticStressCombinedStressStateMxfpFourBSixteenWorstOutputCosine}{0.928\xspace}
% numbers.json:synthetic_stress_combined_stress_state_mxfp4_b16_worst_state_rel_l2
\newcommand{\SyntheticStressCombinedStressStateMxfpFourBSixteenWorstStateRelLTwo}{0.211\xspace}
% numbers.json:synthetic_stress_combined_stress_state_mxfp8_b16_mean_output_cosine
\newcommand{\SyntheticStressCombinedStressStateMxfpEightBSixteenMeanOutputCosine}{0.979\xspace}
% numbers.json:synthetic_stress_combined_stress_state_mxfp8_b16_worst_output_cosine
\newcommand{\SyntheticStressCombinedStressStateMxfpEightBSixteenWorstOutputCosine}{0.935\xspace}
% numbers.json:synthetic_stress_combined_stress_state_mxfp8_b16_worst_state_rel_l2
\newcommand{\SyntheticStressCombinedStressStateMxfpEightBSixteenWorstStateRelLTwo}{0.195\xspace}
% numbers.json:synthetic_stress_outlier_state_mxfp4_b16_mean_output_cosine
\newcommand{\SyntheticStressOutlierStateMxfpFourBSixteenMeanOutputCosine}{0.975\xspace}
% numbers.json:synthetic_stress_outlier_state_mxfp4_b16_worst_output_cosine
\newcommand{\SyntheticStressOutlierStateMxfpFourBSixteenWorstOutputCosine}{0.968\xspace}
% numbers.json:synthetic_stress_outlier_state_mxfp4_b16_worst_state_rel_l2
\newcommand{\SyntheticStressOutlierStateMxfpFourBSixteenWorstStateRelLTwo}{0.159\xspace}
% numbers.json:synthetic_stress_outlier_state_mxfp8_b16_mean_output_cosine
\newcommand{\SyntheticStressOutlierStateMxfpEightBSixteenMeanOutputCosine}{0.981\xspace}
% numbers.json:synthetic_stress_outlier_state_mxfp8_b16_worst_output_cosine
\newcommand{\SyntheticStressOutlierStateMxfpEightBSixteenWorstOutputCosine}{0.973\xspace}
% numbers.json:synthetic_stress_outlier_state_mxfp8_b16_worst_state_rel_l2
\newcommand{\SyntheticStressOutlierStateMxfpEightBSixteenWorstStateRelLTwo}{0.114\xspace}
% numbers.json:usc_baseline_energy_mj
\newcommand{\UscBaselineEnergyMj}{9.5\,mJ\xspace}
% numbers.json:usc_baseline_latency_us
\newcommand{\UscBaselineLatencyUs}{63.2\,$\mu$s\xspace}
% numbers.json:usc_baseline_power_w
\newcommand{\UscBaselinePowerW}{150\,W\xspace}

%
% The generated table paths below assume this file is compiled from the
% repository root or that the paper/ directory is uploaded beside the manuscript.

\section{Results}
\label{sec:results}

\subsection{End-to-End Decode Comparison}

The default MXFP4 persistent-state design reaches
\OursMxfpFourBThirtyTwoLatencyUs per token in RTL cosimulation and uses
\OursMxfpFourBThirtyTwoPowerW of Vivado-estimated post-implementation on-chip
power. Against the cited H100 baseline, this corresponds to
\HeadlineSpeedupVsHOneHundred speedup and
\HeadlineEnergyEfficiencyVsHOneHundred energy efficiency. Against the cited
FPGA baseline, the same design gives \HeadlineSpeedupVsUsc speedup and
\HeadlineEnergyEfficiencyVsUsc energy-efficiency gain relative to the cited
FPGA energy upper bound. These results support the main architectural claim:
the benefit comes from combining native MXFP4 arithmetic with persistent
recurrent-state dataflow, rather than from reduced precision alone.

\begin{table}[t]
\centering
\caption{Batch-one Gated DeltaNet decode comparison. Baseline rows are cited
prior-work values; our row is measured from RTL cosimulation and Vivado
post-implementation estimates.}
\label{tab:results-main}
\begin{tabular}{llll}
\toprule
System & Datapath/state & Latency & Energy \\
\midrule
H100 baseline & HBM state traffic & \HOneHundredBaselineLatencyUs & \HOneHundredBaselineEnergyMj \\
FPGA baseline & Persistent state & \UscBaselineLatencyUs & \UscBaselineEnergyMj \\
Ours & MXFP4 persistent state & \OursMxfpFourBThirtyTwoLatencyUs & \OursMxfpFourBThirtyTwoEnergyMj \\
\bottomrule
\end{tabular}
\end{table}

\subsection{Implementation Evidence}

The implementation closes timing for the Alveo U55C target. HLS C-synthesis
estimates \OursMxfpFourBThirtyTwoCsynthFmaxMhz, and Vivado implementation
reports \OursMxfpFourBThirtyTwoImplWnsNs post-implementation slack. The default
design uses \OursMxfpFourBThirtyTwoLutPct of LUTs,
\OursMxfpFourBThirtyTwoBramPct of BRAM, and \OursMxfpFourBThirtyTwoDspPct of
DSPs. The FP4 element multiply is LUT-based, while the remaining DSP usage
comes from wider HLS-generated fixed-point operations outside the native E2M1
multiply primitive; the source-level native MAC evidence reports no DSP pragma
in the MAC source.

\begin{table}[t]
\centering
\caption{Measured implementation evidence for the default MXFP4 configuration.}
\label{tab:results-implementation}
\begin{tabular}{lp{0.52\linewidth}}
\toprule
Metric & Result \\
\midrule
C-sim parity & \HlsCsimVectorsPassed/\HlsCsimVectorsTotal vectors \\
C/RTL cosim length & \OursMxfpFourBThirtyTwoTokensCosim \\
HLS estimated Fmax & \OursMxfpFourBThirtyTwoCsynthFmaxMhz \\
HLS latency & \OursMxfpFourBThirtyTwoCsynthLatencyCycles \\
RTL cosim latency & \OursMxfpFourBThirtyTwoLatencyCycles \\
Post-implementation slack & \OursMxfpFourBThirtyTwoImplWnsNs \\
Vivado-estimated on-chip power & \OursMxfpFourBThirtyTwoPowerW \\
Resource utilization & \OursMxfpFourBThirtyTwoLutPct LUT, \OursMxfpFourBThirtyTwoBramPct BRAM, \OursMxfpFourBThirtyTwoDspPct DSP \\
\bottomrule
\end{tabular}
\end{table}

\subsection{Parallelism and Block-Size Sweep}

The design sweep shows that the default block-size and parallelism point is the
best current headline configuration. Wider parallelism increases resource use
without improving the measured end-to-end latency in the present artifact,
which indicates that the token path is limited by state movement, interface
packing, and HLS scheduling rather than by the number of raw arithmetic lanes
alone. Smaller parallelism reduces area and power, but it also increases
latency.

\begin{table*}[t]
\centering
\caption{MXFP4 design sweep. Non-default rows are retained as design-space
evidence, not as the headline implementation point.}
\label{tab:results-sweep}
\begin{tabular}{lllll}
\toprule
Configuration & Latency & Power & LUT & DSP \\
\midrule
Default, $B=32$, $P_K=16$, $P_V=8$ &
\OursMxfpFourBThirtyTwoLatencyUs &
\OursMxfpFourBThirtyTwoPowerW &
\OursMxfpFourBThirtyTwoLutPct &
\OursMxfpFourBThirtyTwoDspPct \\
Small fabric, $B=32$, $P_K=8$, $P_V=4$ &
\OursMxfpFourPkEightPvFourBThirtyTwoLatencyUs &
\OursMxfpFourPkEightPvFourBThirtyTwoPowerW &
\OursMxfpFourPkEightPvFourBThirtyTwoLutPct &
\OursMxfpFourPkEightPvFourBThirtyTwoDspPct \\
Wide fabric, $B=32$, $P_K=32$, $P_V=16$ &
\OursMxfpFourPkThirtyTwoPvSixteenBThirtyTwoLatencyUs &
\OursMxfpFourPkThirtyTwoPvSixteenBThirtyTwoPowerW &
\OursMxfpFourPkThirtyTwoPvSixteenBThirtyTwoLutPct &
\OursMxfpFourPkThirtyTwoPvSixteenBThirtyTwoDspPct \\
$B=16$, default parallelism &
\OursMxfpFourBSixteenPkSixteenPvEightLatencyUs &
\OursMxfpFourBSixteenPkSixteenPvEightPowerW &
\OursMxfpFourBSixteenPkSixteenPvEightLutPct &
\OursMxfpFourBSixteenPkSixteenPvEightDspPct \\
$B=16$, wide fabric &
\OursMxfpFourBSixteenPkThirtyTwoPvSixteenLatencyUs &
\OursMxfpFourBSixteenPkThirtyTwoPvSixteenPowerW &
\OursMxfpFourBSixteenPkThirtyTwoPvSixteenLutPct &
\OursMxfpFourBSixteenPkThirtyTwoPvSixteenDspPct \\
\bottomrule
\end{tabular}
\end{table*}

\subsection{Synthetic Numerical Fidelity}

The current accuracy evidence is synthetic-only because the Qwen3-Next capture
status is \QwenCaptureStatus. Under the deterministic synthetic harness, MXFP4
keeps output direction closer to the FP32 reference than the INT4 fallback. With
block size 32, MXFP4 reaches \SyntheticMxfpFourBThirtyTwoOutputCosine output
cosine and \SyntheticMxfpFourBThirtyTwoOutputRelLTwo output relative $L_2$
error. With block size 16, MXFP4 reaches
\SyntheticMxfpFourBSixteenOutputCosine output cosine and
\SyntheticMxfpFourBSixteenOutputRelLTwo output relative $L_2$ error. The INT4
fallback output cosine is \SyntheticIntFourOutputCosine.

\begin{table}[t]
\centering
\caption{Synthetic quantization fidelity. These values should not be described
as Qwen3-Next perplexity or real-model quality results.}
\label{tab:results-accuracy}
\begin{tabular}{lll}
\toprule
Configuration & Output cosine & Output relative $L_2$ \\
\midrule
MXFP4, $B=32$ & \SyntheticMxfpFourBThirtyTwoOutputCosine & \SyntheticMxfpFourBThirtyTwoOutputRelLTwo \\
MXFP4, $B=16$ & \SyntheticMxfpFourBSixteenOutputCosine & \SyntheticMxfpFourBSixteenOutputRelLTwo \\
INT4 fallback & \SyntheticIntFourOutputCosine & -- \\
\bottomrule
\end{tabular}
\end{table}

\subsection{Long-Run State Drift and Traffic Avoidance}

The long-token synthetic stress tests show why recurrent-state precision is a
first-order design choice. Over \SyntheticDriftTokens, MXFP4 state reaches
\SyntheticDriftMxfpFourFinalOutputCosine final output cosine and
\SyntheticDriftMxfpFourFinalStateRelLTwo final state relative $L_2$ error,
while the MXFP8-state fallback reaches
\SyntheticDriftMxfpEightFinalOutputCosine final output cosine and
\SyntheticDriftMxfpEightFinalStateRelLTwo final state relative $L_2$ error.
This supports keeping MXFP4 as the native arithmetic path while retaining
MXFP8 state as the safer fallback for accumulated recurrent-state quality.

\begin{table}[t]
\centering
\caption{Synthetic recurrent-state stress evidence. These tests broaden the
synthetic fidelity boundary when real Qwen3-Next activation capture is
unavailable.}
\label{tab:results-state-stress}
% AUTO-GENERATED -- DO NOT EDIT BY HAND
\begin{tabular}{lrr}
\toprule
Synthetic test & MXFP4 state & MXFP8 state \\
\midrule
1024-token final output cosine
  & 0.643 & 0.933 \\
1024-token final state rel. L2
  & 0.865 & 0.344 \\
256-token boundary output cosine
  & 0.747 & 0.945 \\
Combined-stress worst output cosine
  & 0.928 & 0.935 \\
Outlier-stress worst output cosine
  & 0.968 & 0.973 \\
\bottomrule
\end{tabular}

\end{table}

The storage and traffic accounting also supports the persistent-state dataflow
claim. Including MXFP4 scale metadata, the $B=32$ recurrent state occupies
\StateStorageMxfpFourBThirtyTwoBytes bytes, a
\StateStorageMxfpFourBThirtyTwoReductionVsFpThirtyTwoPct reduction versus
FP32 and a \StateStorageMxfpFourBThirtyTwoReductionVsBfSixteenPct reduction
versus BF16. A naive three-pass off-chip state path would move
\OffchipMxfpFourBThirtyTwoNaiveThreePassBytesPerToken bytes per token for
the same MXFP4 $B=32$ state, whereas the persistent-state datapath has
\OffchipMxfpFourBThirtyTwoPersistentBytesPerToken steady-state off-chip
recurrent-state traffic.

\begin{table}[t]
\centering
\caption{Recurrent-state storage and off-chip traffic accounting for MXFP4
$B=32$.}
\label{tab:results-state-traffic}
% AUTO-GENERATED -- DO NOT EDIT BY HAND
\begin{tabular}{lr}
\toprule
Metric & Value \\
\midrule
MXFP4 B=32 recurrent-state storage & 278528 bytes \\
Reduction vs FP32 & 86.719\% \\
Reduction vs BF16 & 73.438\% \\
Naive three-pass off-chip state traffic & 835584 bytes/token \\
Persistent steady-state off-chip state traffic & 0 bytes/token \\
\bottomrule
\end{tabular}

\end{table}

Overall, the measured results establish a simulation-backed hardware result for
one batch-one Gated DeltaNet decode layer. The strongest supported claims are
RTL-cosim latency, Vivado implementation timing, post-implementation resource
use, and synthetic numerical fidelity. The final manuscript should avoid
claiming completed Qwen3-Next perplexity until realistic activation capture is
available.
